
\documentclass[10pt,landscape,fleqn]{article}
%
\usepackage{pgfplots}
\usepackage[T1]{fontenc}
\usepackage[utf8]{inputenc}
\usepackage[ngerman]{babel}
\usepackage{lmodern}
\usepackage{pgfplots}
\pgfplotsset{compat=1.12}
\usepgfplotslibrary{fillbetween}
\usetikzlibrary{patterns}
%
%\usepackage{pstricks,pst-plot,pst-text,pst-solides3d}
\usepackage{epsfig}
\usepackage{graphicx}
\usepackage[colorlinks=true,
            pdftitle={Mathe II/B},
            pdfauthor={Thomas Carraro},
            pdfpagemode={UseOutlines},pdfstartview={Fit},
            linkcolor=black,urlcolor=blue,
            pdfsubject={Mathematik IV}]{hyperref}
\usepackage{cancel}
%
\usepackage{ifthen}
\newboolean{mitLoes}
%\setboolean{mitLoes}{false}
\setboolean{mitLoes}{true}

\newcounter{sol}
\setcounter{sol}{0}
\ifthenelse{\boolean{mitLoes}}{
\setcounter{sol}{1}
}

\usepackage{amssymb,amsmath,amsthm,am1-aufgaben,bm}
\usepackage{array, hhline}
\usepackage{multirow, multicol}
\usepackage{paralist}
\usepackage{mathtools}

% \usepackage{marvosym,pifont}
% \usepackage{listings}
% \usepackage{algorithmic,algorithm}
% \lstset{language=Matlab,numbers=left}

\newcounter{Aufg}
\setcounter{Aufg}{0}

\newcommand{\Aufgabe}[3][]{
\stepcounter{Aufg}
\subsubsection*{Aufgabe
\arabic{Aufg}}
{#3}
}

\newcommand{\Loesung}[1]{
	\ifthenelse{\boolean{mitLoes}}
	{\subsubsection*{L\"osung:}
		#1}
	{}
	\ifthenelse{\boolean{mitLoes}}{
	\leavevmode
	\newline
	\rule{360pt}{1pt}
\cleardoublepage 
}

}


\textwidth = 270mm
\textheight = 180mm
\leftmargin 5mm
\oddsidemargin = -20mm
\evensidemargin = -20mm
\topmargin = -25mm
\parindent 0cm
\columnsep 2cm

\newcommand{\ds}{\displaystyle}
\newcommand{\rme}{\mathrm{e}}
\newcommand{\rmi}{\mathrm{i}}

\renewcommand{\vec}{\boldsymbol}
\newcommand{\mat}{\boldsymbol}

\newcommand{\ud}{\;\mathrm{d}}
\newcommand{\p}{\partial}
\newcommand{\dO}{\d \vec o}
\newcommand{\dA}{\d \vec A}
\newcommand{\dX}{\d \vec x}

\renewcommand{\div}{\operatorname{div}}
\newcommand{\rot}{\operatorname{rot}}
\newcommand{\grad}{\operatorname{grad}}

\newcommand{\skalar}[1]{\left\langle{#1}\right\rangle}

\renewcommand{\d}{\,{\rm d}}
\newcommand{\x}{\ensuremath{{\vec x}}}
\newcommand{\ve}[3]{\begin{pmatrix} #1\\ #2\\ #3\end{pmatrix}}
\newcommand{\e}{\ensuremath{{\bf e}}}

\begin{document}



\begin{twocolumn}



\href{http://mb-mathe.hsu-hh.de/bause}
{\large\bf Helmut-Schmidt-Universit\"at}\\
\href{http://mb-mathe.hsu-hh.de/bause}
{\large\bf Universit\"at der Bundeswehr Hamburg}\\
\href{http://mb-mathe.hsu-hh.de/bause}
{\large\bf Fakult\"at f\"ur Maschinenbau und\,Bauingenieurwesen}\\[1ex]
\href{http://mb-mathe.hsu-hh.de/bause}
{\large\bf Prof. Dr.\ Thomas Carraro}\\
\\[-28mm]
\phantom{.}
\hfill\includegraphics[height=20mm]{HSU_RGB}

\vspace*{8mm}
{\Large\bf Mathematik II/B (WI/ET) \hfill Zusatzblatt}\\[1ex]
{\Large WT 2025 \hfill  Taylor-Entwicklung}


\vspace*{1mm}
\rule{360pt}{2pt}

\vspace*{0ex}
% \paragraph*{Ex1}
%
\Aufgabe{}{
Gegeben sei die Funktion $$
f(x) = \frac{1}{(x+1)^2},
$$
\begin{abc}
\item Bestimmen Sie das Taylor-Polynom 3. Ordnung $T_3(x)$ von $f (x)$ um den
	  Punkt $x_0 = 0$.
\item Geben Sie das Restglied $R_3(x;0)$ an. 
\item Schätzen Sie das Restglied für $|x| < \frac{1}{2}$ ab.
\end{abc}
\textbf{Hinweis:} Es ist nicht notwendig das Restglied optimal abzuschätzen.
}


\Loesung{
\begin{abc}
\item Die Taylor-Entwicklung von \( f(x) \) um \( x_0 = 0 \) ist:
\[
f(x) = f(0) + f'(0)x + \frac{f''(0)}{2!}x^2 + \frac{f'''(0)}{3!}x^3 \cdots
\]

\begin{align*}
f(0)  &= 1 \\
f'(x) &= -\frac{2}{(x+1)^3} \\
f'(0) &= -2 \\
f''(x) &= \frac{6}{(x+1)^4} \\
f''(0) &= 6 \\
f'''(x) &= -\frac{24}{(x+1)^5} \\
f'''(0) &= -24
\end{align*}


Damit ergibt sich das Taylor-Polynom zu
\[
T_3(x) = f(0) + f'(0)x + \frac{f''(0)}{2!}x^2 +  \frac{f'''(0)}{3!}x^3 \\
= 1-2x+3x^2-4x^3.
\]

\item Das Restglied wird wie folgt bestimmt:
\[
R_3(x; 0) = \frac{f^{(4)}(\xi)}{4!}x^4,
\]
wobei $-\frac{1}{2}\leq \xi \leq \frac{1}{2}$ gilt.

\[
f^{(4)}(x) = \frac{120}{(x+1)^6}.
\]

Das Restglied wird damit 
\[
R_3(x; 0) = \frac{f^{(4)}(\xi)}{4!}x^4 = \frac{120}{24(\xi+1)^6}x^4.
\]

\item For \( |x| < \frac{1}{2} \):
\[
|R_3(x; 0)| = \left|\frac{120}{24(\xi+1)^6}x^4\right| = \frac{5}{(\xi+1)^6}x^4
\leq \frac{5}{(\frac{1}{2})^6} \left( \frac{1}{2}\right) ^4 = \frac{5 \cdot 64}{16} = 20.
\]

D.h.
\[
|R_3(x; 0)| \leq 20.
\]

	  
\end{abc}

}

\Aufgabe{}{
	Gegeben sei die Funktion 
	$$
	f(x) = e^{-2x}.
	$$
	\begin{abc}
		\item Bestimmen Sie das Taylor-Polynom 3. Ordnung \( T_3(x) \) von \( f(x) \) um den Punkt \( x_0 = 0 \).
		\item Geben Sie das Restglied \( R_3(x;0) \) an. 
		\item Schätzen Sie das Restglied für \( |x| < \frac{1}{2} \) ab.
	\end{abc}
}

\Loesung{
	\begin{abc}
		\item Die Taylor-Entwicklung von \( f(x) \) um \( x_0 = 0 \) ist:
		\[
		T_3(x) = f(0) + f'(0)x + \frac{f''(0)}{2!}x^2 + \frac{f'''(0)}{3!}x^3 + \dots
		\]
		
		Berechnung der Ableitungen:
		
		\begin{align*}
			f(0) &= 1 \\
			f'(x) &= -2e^{-2x} \quad \Rightarrow \quad f'(0) = -2 \\
			f''(x) &= 4e^{-2x} \quad \Rightarrow \quad f''(0) = 4 \\
			f'''(x) &= -8e^{-2x} \quad \Rightarrow \quad f'''(0) = -8
		\end{align*}
		
		Damit ergibt sich das Taylor-Polynom dritter Ordnung:
		
		\[
		T_3(x) = f(0) + f'(0)x + \frac{f''(0)}{2!}x^2 + \frac{f'''(0)}{3!}x^3
		\]
		
		\[
		T_3(x) = 1 - 2x + 2x^2 - \frac{4}{3}x^3.
		\]
		
		\item Das Restglied wird durch die Lagrange-Formel gegeben:
		
		\[
		R_3(x; 0) = \frac{f^{(4)}(\xi)}{4!}x^4,
		\]
		
		
		
		Die vierte Ableitung von \( f(x) \) ist:
		
		\[
		f^{(4)}(x) = 16e^{-2x}.
		\]
		
		Damit folgt für das Restglied:
		
		\[
		R_3(x; 0) = \frac{16e^{-2\xi}}{24}x^4 = \frac{2}{3}e^{-2\xi} x^4.
		\]
		
		\item Für \( |x| < \frac{1}{2} \) setzen wir die obere Schranke für \( e^{-2\xi} \):
		
		Da \( e^{-2\xi} \) auf \( [-\frac{1}{2}, \frac{1}{2}] \) maximal ist für \( \xi = -\frac{1}{2} \), gilt:
		
		\[
		e^{-2\xi} \leq e.
		\]
		
		Damit ergibt sich die Abschätzung:
		
		\[
		|R_3(x; 0)| \leq \frac{2}{3} e \cdot \left(\frac{1}{2}\right)^4.
		\]
		
		\[
		= \frac{2}{3} e \cdot \frac{1}{16} = \frac{2e}{48} = \frac{e}{24}.
		\]
		
		Also:
		
		\[
		|R_3(x; 0)| \leq \frac{e}{24}.
		\]
		
	\end{abc}
}

% 



%%%%%%%%%%%%%%%%%%%%%%%%%%%%%%%%%%%%%%%%%%%%%%%%%%%%%%%%%%%%%%%%%%%%%%%%%%%%%%%%


%%%%%%%%%%%%%%%%%%%%%%%%%%%%%%%%%%%%%%%%%%%%%%%%%%%%%%%%%%%%%%%%%%%%%%%%%%%%%%%%
\end{twocolumn}

\end{document}

%%%%%%%%%%%%%%%%%%%%%%%%%%%%%%%%%%%%%%%%%%%%%%%%%%%%%%%%%%%%%%%%%%%%%%%%%%%%%%%%
